\section{中央 pencil の Rees 模型の基底変換とアフィン半直交分解}
\label{sec:central-basechange}
本節では第\ref{sec:restored-rees-sod}節に残された問題を、
中央二変数環に関する\emph{アフィン錐の段階}で解く。
次の議論は与えられた中央正則列について完全に代数的であり、
特定の量子化や形式的な変形理論には依存しない。
ただし最後に述べるように、射影圏への降下は別の問題である。

\begin{theorem}[中央正則列に関する次数付き自由性]
\label{thm:central-graded-freeness}
$k$ を体とし、$A=\bigoplus_{n\ge0}A_n$ を
$A_0=k$、$\dim_kA_n<\infty$ の連結次数付き結合代数とする。
$s,t\in Z(A)_1$ は中央元で、$s$ は $A$ の非零因子、
$t$ は $A/sA$ の非零因子であるとする。
$C=k[s,t]$、$D=A/(s,t)$ と置くと、
\begin{enumerate}
\item 自然な写像 $k[S,T]\to A$、$S\mapsto s,T\mapsto t$ は単射である。
\item $D$ の任意の斉次 $k$-基底を $A$ に斉次に持ち上げると、
それらの持ち上げは $C$-加群として $A$ の自由基底をなす。
\item 特に $A$ は $C$ 上忠実平坦で、Hilbert 級数は
\[
 H_A(z)=\frac{H_D(z)}{(1-z)^2}
\]
を満たす。
\end{enumerate}
\end{theorem}
\begin{proof}
$s,t$ の正則性から
$0\to A(-1)\xrightarrow{s}A\to A/sA\to0$ と
$0\to (A/sA)(-1)\xrightarrow{t}A/sA\to D\to0$
は完全である。従って第3項の Hilbert 級数の公式を得る。
$D$ の各次数の基底を $A$ に持ち上げた集合を $\{b_j\}$ とする。
非負次数について帰納すれば、$A_n$ の任意の元は $b_j$ の
$C$-線形結合で表せる。実際、$D_n$ の基底を引き去ると
$(s,t)A$ に属するので、一つ低い次数へ帰着する。
従って
$F=\bigoplus_jC(-\deg b_j)\twoheadrightarrow A$
は次数付き全射である。
$H_F=H_D/(1-z)^2=H_A$ なので、各次数で全射は同じ有限次元
$k$-ベクトル空間の間の同型である。従って $F\simeq A$。
もし $P(s,t)=0$ という非零の多項式関係があれば、
得られた自由 $k[S,T]$-加群構造に反する。
$D_0=k$ の持ち上げは $1$ と取れるので、自由加群は零でなく、
忠実平坦性も従う。
\end{proof}

\begin{theorem}[Rees 環の可換 blow-up からの基底変換]
\label{thm:central-rees-basechange}
定理\ref{thm:central-graded-freeness}の仮定の下で $I=(s,t)A$ とする。
$m=(s,t)\subset C$、$\mathcal R_C(m)=\bigoplus_{r\ge0}m^r u^r$ とおけば
\[
\boxed{\quad
 \mathcal R_A(I)\;\simeq\; A\otimes_C\mathcal R_C(m)
 \;\simeq\; A[U,V]/(tU-sV).
 \quad}
\]
ここで $U,V$ は $A$ と可換で、Rees 次数は $1$、元の次数も $1$ である。
$Y=\operatorname{Bl}_0\mathbb A^2=\operatorname{Proj}_{C}\mathcal R_C(m)$
とし、$f:Y\to\operatorname{Spec}C$ と置くと、
Rees 方向の非可換相対 $\operatorname{Proj}$ は
$Y$ 上の準連接な代数の層
\[
  \mathscr A=\mathcal O_Y\otimes_C A
\]
で記述できる。その二つの標準アフィンチャートは
\[
\mathscr A(U\ne0)\simeq A[w]/(t-sw),\qquad
\mathscr A(V\ne0)\simeq A[v]/(s-tv),
\]
ここで $w=V/U$、$v=U/V$ は中央元で、交叉では $wv=1$ である。
例外因子 $E=f^{-1}(0)\simeq\mathbb P^1$ に沿う制限は
\[
  \mathscr A|_E\simeq D\otimes_k\mathcal O_{\mathbb P^1}.
\]
さらに元の次数を $z$、Rees 次数を $w$ で数えた二変数 Hilbert 級数は
\[
 H_{\mathcal R_A(I)}(z,w)
   =H_A(z)\frac{1-z^2w}{(1-zw)^2}.
\]
\end{theorem}
\begin{proof}
$I^r=m^rA$ であり、$A$ は $C$-平坦だから
$\mathcal R_C(m)\hookrightarrow C[u]$ に
$A\otimes_C-$ を施した写像も単射である。その像は
$\bigoplus_{r\ge0}I^ru^r$ である。
可換二変数の blow-up に対する標準表示
$\mathcal R_C(m)\simeq C[U,V]/(tU-sV)$ を
基底変換すれば最初の二つの同型を得る。
二つのチャートは $U$ または $V$ で局所化し、
Rees 次数零部分を取ることで得られる。
$s=t=0$ を課すと $D[w]$ と $D[v]$ が生じる。
また二重次数 $(\deg_A,\deg_{\mathrm{Rees}})$ は
$\deg U=\deg V=(1,1)$、$\deg(tU-sV)=(2,1)$
なので、正則な関係式一つを除すことから Hilbert 級数の式が従う。
同じ式は $I^r=m^rA$ と $C$-自由性からも直接得られる。
\end{proof}

\begin{theorem}[アフィン非可換 blow-up の Orlov 分解]
\label{thm:central-affine-orlov}
さらに $A$ を左右 noetherian とする。
$\mathscr A=\mathcal O_Y\otimes_C A$ 上の perfect な右加群複体の
dg 圏を $\operatorname{Perf}(Y,\mathscr A)$ と書く。
$\iota:E\hookrightarrow Y$ を例外因子とし、
$\operatorname{Perf}(D)$ の対象 $M$ に
\[
 \Psi(M)=\iota_*\bigl((M\otimes_k\mathcal O_E)\otimes\mathcal O_E(-1)\bigr),
 \qquad
 \Phi(N)=N\otimes_A^{\mathbf L}\mathscr A
 \quad(N\in\operatorname{Perf}(A))
\]
を対応させる。両関手は dg fully faithful であり、
\[
\boxed{\quad
\operatorname{Perf}(Y,\mathscr A)
 =\langle\Psi\operatorname{Perf}(D),\,
               \Phi\operatorname{Perf}(A)\rangle.
\quad}
\]
混合 Hom は
\[
 \operatorname{RHom}_{Y,\mathscr A}(\Phi N,\Psi M)=0,
 \qquad
 \operatorname{RHom}_{Y,\mathscr A}(\Psi M,\Phi N)
 \simeq\operatorname{RHom}_{D}
       (M,N\otimes_A^{\mathbf L}D)[-1]
\]
で与えられる。この SOD は第\ref{thm:restored-categorical-gluing}定理の
上三角 bimodule の具体的な\emph{アフィン}実現である。
特にRees と gluing の比較は、
\emph{アフィン錐については}成立する。
\end{theorem}
\begin{proof}
可換な $f:Y\to\operatorname{Spec}C$ は滑らかな点の余次元二 blow-up
なので、Orlov の定理\cite{OrlovBlowup}は
\[
 \operatorname{Perf}(Y)
 =\langle\iota_*\mathcal O_E(-1)\otimes\operatorname{Perf}(k),
                 f^*\operatorname{Perf}(C)\rangle
\]
という $\operatorname{Perf}(C)$-線形の分解を与える。
$A$ の $C$-平坦性によって、これを $A$-係数の加群圏へ
基底変換できる。具体的には、$Y$ の二つのアフィンチャート上で
$\operatorname{Mod}(\mathcal O_Y)\otimes_{\operatorname{Mod}(C)}
\operatorname{Mod}(A)$ は
$\operatorname{Mod}(\mathcal O_Y\otimes_C A)$ と同値であり、
この同値はチャートの交わりと両立する。
従って準連接導来加群圏を貼り合わせ、compact objects を取れば、
$\operatorname{Perf}(Y)$ の上記分解を
$\operatorname{Perf}(Y,\mathscr A)$ へ移せる。
この基底変換の完備化・compact 性は
可換側の Orlov 分解の局所的構成からも確認できる
（半直交分解の基底変換について\cite{KuznetsovSODBaseChange}も参照）。

直接、完全忠実性と混合 Hom を確認しておく。
$Rf_*\mathcal O_Y=C$、$R^if_*\mathcal O_Y=0\;(i>0)$ と
$C$-平坦性により $Rf_*\mathscr A=A$ となり、$\Phi$ は fully faithful。
$E$ は Cartier 因子で、その法束は $\mathcal O_E(-1)$ である。
$\iota^!\Phi N\simeq
 ((N\otimes_A^{\mathbf L}D)\otimes_k\mathcal O_E(-1))[-1]$
から第二の混合 Hom を得る。
第一の混合 Hom は $R\Gamma(E,\mathcal O_E(-1))=0$ から消える。
同じ Cartier resolution を $\Psi M,\Psi M'$ に使うと
$\operatorname{RHom}(\Psi M,\Psi M')\simeq
 \operatorname{RHom}_D(M,M')$ となる。
全生成性は上記の $C$-線形 Orlov 分解の基底変換による。
\end{proof}

\begin{proposition}[射影圏への降下の際に残る非正則頂点]
\label{prop:central-vertex-obstruction}
最大 pencil の中央商が
$D=B(E,L,\tau)$、$E$ が楕円曲線で $\deg L=6$ であるとする。
このとき $A$ と $D$ の連結次数付き代数としての自明加群 $k$ は
いずれも無限射影次元を持つ。実際
\[
 H_D(z)=\frac{1+4z+z^2}{(1-z)^2},\qquad
 H_A(z)=\frac{1+4z+z^2}{(1-z)^4}
   =\sum_{n\ge0}(n+1)^3z^n.
\]
従って定理\ref{thm:central-affine-orlov}の
アフィン perfect 圏の SOD を、
単に「有限次元加群で Verdier 商を取る」ことによって
$\operatorname{Perf}(\operatorname{qgr}A)$ の SOD にすることはできない。
これはそのような射影 SOD が存在しないという主張ではなく、
\emph{導来 torsion 局所化を別途構成する必要がある}という障害である。
\end{proposition}
\begin{proof}
$\deg L=6$ の楕円曲線では $h^0(E,L_n)=6n$（$n>0$）なので
表示した $H_D$ を得る。中央正則列から $H_A=H_D/(1-z)^2$。
もし連結次数付き右 noetherian 代数 $R=A$ または $D$ 上で
$k=R/R_{\ge1}$ の射影次元が有限なら、有限個の有限生成次数付き
自由加群からなる $k$ の分解があり、Euler--Poincar\'e の公式から
$1/H_R(z)$ は Laurent 多項式となる。
しかし $1+4z+z^2$ は $(1-z)$ と互いに素であるため
\[
 H_D(z)^{-1}=\frac{(1-z)^2}{1+4z+z^2},\qquad
 H_A(z)^{-1}=\frac{(1-z)^4}{1+4z+z^2}
\]
はいずれも Laurent 多項式にならない。矛盾する。
従って irrelevant torsion の基本対象 $k$ は
$\operatorname{Perf}^{\mathrm{gr}}(A)$ および
$\operatorname{Perf}^{\mathrm{gr}}(D)$ に含まれない。
射影化ではこの非 perfect torsion も消すので、
アフィン perfect 圏の単純な商だけで射影圏を記述することはできない。
\end{proof}

\begin{remark}[次の具体的な比較]
上記の証明で Rees の\emph{アフィン} gluing は決定した。
射影三次元模型に対しては、元の次数と Rees 次数の\emph{二重次数付き}
導来加群圏を取り、まず Rees 方向の相対 Proj を構成し、
次に元の irrelevant torsion を消す操作が必要である。
そのとき投影関手が torsion 局所化と両立するか、
$\operatorname{qgr}D\simeq\operatorname{coh}E$ の例外成分が
正しい dg 埋込みとして残るか、を検証しなければならない。
次の目標は二重次数付き局所化後の
$\langle\operatorname{Perf}(E),\operatorname{Perf}(\operatorname{qgr}A)\rangle$
という SOD の成立を実際に証明することであり、
本節のアフィン結果から無条件に結論しない。
\end{remark}

\begin{proposition}[可換二点 blow-up の反標準性の境界]
\label{prop:anticanonical-commutative-boundary}
$Y$ を $\mathbb P^2$ の二点での逐次 blow-up とする。
二点が元の平面上で相異なれば $-K_Y$ は ample で、
$\operatorname{coh}Y\simeq\operatorname{qgr}R(Y,-K_Y)$ である。
二点目が最初の例外曲線上にあれば、$-K_Y$ は big, nef だが
ample ではなく、最初の例外曲線の狭義変換を収縮する。
両場合で反標準環の Hilbert 関数は同じく
\[
 h^0(Y,-nK_Y)=1+\frac{7n(n+1)}2\quad(n\ge0)
\]
である。従ってこの Hilbert 関数の一致は、
反標準 section functor の ample 性を保証しない。
\end{proposition}
\begin{proof}
元の平面上で二点が相異なる場合、Mori 錐は二つの例外曲線
$E_1,E_2$ と二点を通る直線の狭義変換
$L=H-E_1-E_2$ により生成される。
$-K_Y=3H-E_1-E_2$ はこれらと交点数 $1$ を持ち、
$(-K_Y)^2=7$ なので Nakai--Moishezon により ample。
従って通常の Serre の定理により反標準環から coherent 圏を復元できる。
二点目が infinitely near なら $C=E_1-E_2$ は
$C^2=-2$、$(-K_Y)\cdot C=0$ を満たす。
この二点配置の反標準類は big, nef で、$C$ を収縮するから
ample ではない（第\ref{ex:restored-weak-dp}注意も参照）。
両場合 $K_Y^2=7$, $\chi(\mathcal O_Y)=1$ なので
Riemann--Roch により $\chi(-nK_Y)=1+7n(n+1)/2$。
$n>0$ の高次コホモロジーは
$-nK_Y=K_Y+(n+1)(-K_Y)$ への Kawamata--Viehweg 消滅から零であり、
$n=0$ も有理曲面のコホモロジーから従う。
\end{proof}
